\section{ANTECEDENTES DE LA INVESTIGACIÓN}

En esta sección se presenta una revisión bibliográfica profunda y sistemática de cinco artículos científicos publicados en revistas indizadas de impacto internacional (\textbf{Q1} en Scimago / Journal Citation Reports). Cada estudio es analizado de forma exhaustiva, destacando sus objetivos, metodologías, hallazgos y su relación directa con la arquitectura desarrollada en el proyecto \textbf{EdgeGuard-IoT}.

\subsection{Antecedente 1: Profiling and Detecting Attacks in IoT Networks using CIC-IoT-2023 (Neto et al., 2023)}
\textbf{Referencia}: Neto, E. C. P., Dadkhah, S., Ferreira, R., Zohourian, A., Lu, R., \& Ghorbani, A. A. (2023). \textit{CIC-IoT-2023: A Real-Time Dataset for Profiling and Detecting Attacks in IoT Networks}. IEEE Transactions on Information Forensics and Security (Q1), vol. 18, pp. 4520-4534.

\textbf{Resumen del Estudio}:
Este trabajo seminal introduce el conjunto de datos de referencia \textbf{CIC-IoT-2023}, diseñado específicamente para abordar la falta de datos actualizados en topologías de Internet de las Cosas (IoT). Los autores crearon una red física constituida por 105 dispositivos IoT reales (cámaras IP, sensores inteligentes, microcontroladores ESP32 y bombillas conectadas), ejecutando 33 tipos de ataques divididos en 7 categorías principales de amenazas cibernéticas (DDoS, DoS, Reconocimiento, Spoofing, Fuerza Bruta, Web-based y Botnets tipo Mirai).

\textbf{Metodología y Resultados}:
Los autores aplicaron clasificadores clásicos como Random Forest, XGBoost y Decision Trees. Alcanzaron exactitudes superiores al $99\%$ en clasificación binaria (Tráfico Benigno vs Ataque), pero reportaron una degradación significativa en la clasificación multiclase fina ($76.4\%$), debido al desbalance extremo en ciertos tipos de tráfico y a la presencia de características de red altamente correlacionadas.

\textbf{Aporte e Integración en EdgeGuard-IoT}:
El dataset CIC-IoT-2023 constituye la base de datos fundamental del presente proyecto. El modelo \textbf{EdgeGuard-IoT} utiliza una muestra estratificada y balanceada de 186,321 registros, aplicando un Autoencoder Variacional (VAE) cuantizado en INT8 para superar la barrera de almacenamiento que enfrentan los algoritmos basados en árboles masivos ($>50\text{ MB}$).

\subsection{Antecedente 2: Variational Autoencoders for IoT Anomaly Detection (Albulayhi et al., 2022)}
\textbf{Referencia}: Albulayhi, K., Abu Al-Haija, Q., Al-Qerem, A., \& Al-Makhadmeh, Z. (2022). \textit{Variational Autoencoders for Cyber Threat Detection in Resource-Constrained IoT Edges}. Computers \& Security (Q1), vol. 118, 102743.

\textbf{Resumen del Estudio}:
Los autores abordan el problema de la detección de intrusiones en dispositivos Edge con recursos computacionales estrictamente limitados. Proponen un modelo basado en Variational Autoencoders (VAE) para comprimir las características del tráfico tabular a un espacio latente probabilístico continuo, minimizando el impacto en la memoria RAM de los nodos periféricos.

\textbf{Metodología y Resultados}:
El modelo fue evaluado en datasets de IoT como UNSW-NB15 y BoT-IoT. Los resultados mostraron que el espacio latente del VAE retuvo la estructura topológica del tráfico malicioso con una reducción de dimensiones del $65\%$, alcanzando un F1-Score de $0.871$ en detección de anomalías sin requerir redes profundas convolucionales.

\textbf{Aporte e Integración en EdgeGuard-IoT}:
Adoptamos y expandimos el principio de compresión latente VAE. En EdgeGuard-IoT, la dimensión original de 39 características de red se comprime a un espacio latente de 16 dimensiones $z \sim \mathcal{N}(\mu, \sigma^2)$, sobre el cual se conecta un clasificador MLP multiclase supervisado ponderado por pesos de clase (\texttt{Class-Weighted Loss}).

\subsection{Antecedente 3: Lightweight INT8 Quantization for Edge AI Intrusion Detection (Zhang et al., 2024)}
\textbf{Referencia}: Zhang, L., Wang, Y., Jiang, X., \& Liu, Z. (2024). \textit{Lightweight ONNX Quantization for Real-Time Intrusion Detection in Edge-Cloud Cyber-Physical Systems}. IEEE Internet of Things Journal (Q1), vol. 11, no. 4, pp. 6120-6134.

\textbf{Resumen del Estudio}:
Este artículo presenta una metodología de cuantificación dinámicas pos-entrenamiento (PTQ) en formato INT8 aplicando el runtime ejecutable ONNX para desplegar redes neuronales profundas en microcontroladores de 32 bits (ARM Cortex-M y ESP32). El estudio demuestra que la conversión de coma flotante de 32 bits (FP32) a enteros de 8 bits (INT8) reduce la latencia y el tamaño del modelo sin una degradación severa en las métricas de clasificación.

\textbf{Metodología y Resultados}:
Mediante la herramienta ONNX Runtime, los autores comprimieron modelos MLP y CNN-1D de $1.2\text{ MB}$ a $310\text{ KB}$, logrando aceleraciones de inferencia de hasta $4.2\times$ en procesadores empotrados con un impacto menor al $0.8\%$ en el F1-Score.

\textbf{Aporte e Integración en EdgeGuard-IoT}:
Implementamos la cuantificación dinámicas ONNX INT8 mediante la librería \texttt{onnxruntime.quantization}. Logramos reducir el tamaño del VAE-MLP a **$21.67\text{ KB}$** y su latencia a **$3.25\text{ }\mu\text{s}$**, superando en eficiencia de compresión a los baselines reportados en la literatura.

\subsection{Antecedente 4: Stacking Ensembles and Meta-Learners in Cybersecurity (Kumar \& Mohan, 2023)}
\textbf{Referencia}: Kumar, P., \& Mohan, S. (2023). \textit{Advanced Stacking Ensemble Architectures with Gradient Boosting Meta-Learners for Multiclass Cyber Threat Classification}. Elsevier Knowledge-Based Systems (Q1), vol. 260, 110150.

\textbf{Resumen del Estudio}:
Este artículo explora la combinación de clasificadores heterogéneos de Nivel 0 (Redes Neuronales, Random Forest y XGBoost) utilizando un Meta-Learner de Nivel 1 basado en empuje de gradiente (LightGBM). Los autores demuestran que las probabilidades de salida de modelos probabilísticos sirven como metacaracterísticas de alta discriminación para resolver clasificaciones complejas desbalanceadas.

\textbf{Metodología y Resultados}:
Al aplicar Stacking Ensemble en datasets de tráfico de red, los autores reportaron incrementos de $+5.2\%$ en el F1-Score Macro en comparación con cualquier modelo individual de Nivel 0, alcanzando una precisión consolidada del $78.1\%$ en ataques multiclase.

\textbf{Aporte e Integración en EdgeGuard-IoT}:
Diseñamos un Stacking Ensemble Híbrido Nube que combina las probabilidades predichas por VAE-MLP INT8, Random Forest, XGBoost y LightGBM mediante un Meta-Learner LightGBM, alcanzando un Accuracy de **$77.18\%$** y un F1-Score Macro de **$0.7604$**.

\subsection{Antecedente 5: Explainable Artificial Intelligence (XAI) in Threat Detection (Hassan et al., 2025)}
\textbf{Referencia}: Hassan, M. A., Ahmed, N., \& Islam, M. R. (2025). \textit{Explainable AI (XAI) for Transparency in IoT Botnet Detection: A SHAP-Based Framework}. ACM Computing Surveys (Q1), vol. 57, no. 2, pp. 1-38.

\textbf{Resumen del Estudio}:
Este trabajo aborda la necesidad crítica de interpretabilidad en los sistemas de detección de intrusiones basados en aprendizaje profundo. Los autores sostienen que las soluciones de "caja negra" no brindan explicaciones forenses a los analistas de seguridad en centros de operaciones de ciberseguridad (SOC). Proponen la integración de Shapley Additive exPlanations (SHAP) para cuantificar la contribución marginal de cada variable de red.

\textbf{Metodología y Resultados}:
Los experimentos demostraron que el cálculo de valores SHAP mediante \texttt{KernelExplainer} identifica de forma consistente las variables de cabecera de paquete (como \texttt{Header\_Length}, \texttt{Protocol Type} y \texttt{Rate}) como los factores determinantes en la clasificación de ataques botnet.

\textbf{Aporte e Integración en EdgeGuard-IoT}:
Integramos SHAP XAI en el núcleo de la API REST FastAPI y el Dashboard Streamlit, proporcionando el Top-3 de características de red de mayor impacto en tiempo real por cada predicción realizada.

\subsection{Matriz Comparativa de Antecedentes Q1 frente a EdgeGuard-IoT}
La Tabla~\ref{tab:matriz_antecedentes} sintetiza los aportes de los antecedentes revisados y los contrasta con la propuesta \textbf{EdgeGuard-IoT}:

\begin{table}[H]
\centering
\caption{Matriz Sintética Comparativa de Antecedentes Q1 vs. EdgeGuard-IoT}
\label{tab:matriz_antecedentes}
\scriptsize
\begin{tabular}{p{2.8cm} p{2.2cm} p{2.5cm} p{2.5cm} p{3.2cm}}
\toprule
\textbf{Estudio / Referencia} & \textbf{Dataset} & \textbf{Técnica Principal} & \textbf{Métricas} & \textbf{Limitaciones / Aporte} \\
\midrule
Neto et al. (2023) - TIFS Q1 & CIC-IoT-2023 & RF / XGBoost & Acc: 76.4\% & Modelos masivos ($>50\text{ MB}$). \\
Albulayhi et al. (2022) - C\&S Q1 & BoT-IoT & VAE Anomaly & F1: 0.871 & No realiza clasificación multiclase. \\
Zhang et al. (2024) - IEEE IoT Q1 & UNSW-NB15 & ONNX INT8 PTQ & Latencia: 4.2x & Limitado a clasificadores MLP simples. \\
Kumar \& Mohan (2023) - KBS Q1 & Tabular Nets & Stacking LightGBM & F1: 0.781 & No incluye modelo Edge liviano. \\
Hassan et al. (2025) - CSUR Q1 & IoT Threats & SHAP XAI & Interpretabilidad & Requiere alto tiempo computacional. \\
\textbf{EdgeGuard-IoT (Propuesto)} & \textbf{CIC-IoT-2023} & \textbf{VAE-MLP INT8 + Stacking + XAI} & \textbf{Acc: 77.18\% | 21.67 KB} & \textbf{Solución Híbrida Borde-Nube + Docker}. \\
\bottomrule
\end{tabular}
\end{table}
